\section{Methods}
\label{sec:methods}

\subsection{Stage I: confirmatory audit and evidence boundary}

The diagnostic stage was organized as a frozen confirmatory benchmark rather than a post hoc search for favorable endpoint--method combinations. Development splits and seed~99 were used only for implementation checks and exploratory debugging. Scientific conclusions were restricted to previously unseen confirmatory splits: seeds 101--110 for random and scaffold designs and seeds 101--105 for similarity-cluster designs. Model families, imbalance regimes, conformal methods, applicability-domain definitions, selective-prediction metrics, and stopping rules were fixed before confirmatory outputs were inspected. No endpoint, model, seed, threshold, or unfavorable result was removed after inspection. This separation follows recommendations that molecular machine-learning comparisons distinguish practical effects, paired evaluation units, and repeated-split uncertainty rather than treating model architectures as independent replicates \citep{ash2025comparison}.

\subsection{Endpoint acquisition and structure standardization}

Four public MoleculeNet/DeepChem endpoints were selected to provide two complementary classification and regression stress-test pairs \citep{wu2018moleculenet}. BBBP used the binary \texttt{p\_np} endpoint, ClinTox used \texttt{CT\_TOX}, ESOL used measured log aqueous solubility in mol\,L$^{-1}$, and Lipophilicity used the experimental \texttt{exp} target. The panel spans different sample sizes, class prevalences, and continuous-target distributions, but was not intended to represent the full diversity of ADMET endpoints. BBBP is historically associated with in silico blood--brain barrier penetration modeling \citep{martins2012bbbp}, whereas ESOL originates from the Delaney aqueous-solubility collection \citep{delaney2004esol}.

Raw files were downloaded from the public DeepChem dataset distribution and hashed with SHA-256. Rows with missing structures or targets were removed, and SMILES strings were canonicalized with RDKit. Molecules that could not be parsed after canonicalization were discarded. Duplicate canonical structures were resolved before featurization. For classification, the majority label among duplicate records was retained; an exact tie between conflicting labels was treated as ambiguous and removed. For regression, duplicate target values were aggregated by their median. Table~\ref{tab:endpoints} reports the resulting core datasets. The ClinTox positive prevalence of 6.45\% was retained rather than artificially balanced outside the training-only resampling regime.

\begin{table}[!htbp]
\centering
\caption{\textbf{Core endpoint composition after structure cleaning and duplicate resolution.} Classification prevalence is the fraction labeled positive. Regression ranges are reported in the native target scale distributed with MoleculeNet/DeepChem.}
\label{tab:endpoints}
\small
\resizebox{\textwidth}{!}{%
\begin{tabular}{llrrrrl}
\toprule
Endpoint & Task & Raw $n$ & Clean $n$ & Removed & Conflicts & Prevalence or target range \\
\midrule
BBBP & Classification & 2050 & 1965 & 85 & 10 & 0.7634 positive \\
ClinTox & Classification & 1484 & 1442 & 42 & 19 & 0.0645 positive \\
ESOL & Regression & 1128 & 1117 & 11 & 6 & $-11.60$ to $1.58$ \\
Lipophilicity & Regression & 4200 & 4200 & 0 & 0 & $-1.50$ to $4.50$ \\
\bottomrule
\end{tabular}
}
\end{table}

\subsection{Molecular representation}

Canonical structures were encoded as 2048-bit binary Morgan fingerprints with radius~2, corresponding to ECFP4 \citep{rogers2010ecfp}. The same frozen fingerprint matrix was used for model fitting, similarity-based splitting, density-ratio estimation, and applicability-domain analysis. This representation-controlled design was chosen to isolate reliability diagnostics from changes in learned representation and was not intended as an architecture leaderboard. Modern graph, geometry-aware, and foundation-model encoders may change predictive performance and uncertainty behavior, but varying the representation within this study would confound whether a difference arose from the diagnostic or from the molecular feature space itself \citep{heid2024chemprop,fang2022geometry,kim2024quantum,zhao2026revisiting}.

\subsection{Random, scaffold, and similarity-cluster splits}

Every confirmatory split assigned approximately 60\% of molecules to training, 20\% to calibration, and 20\% to testing. Models and training-only resampling used the training role. Calibration, conformal thresholds, density-ratio estimation, and learned error-scale transformations used only training and/or calibration information as specified below. Test labels were never used for model fitting, split construction, threshold selection, or method selection.

Random classification splits were generated in two stages with stratification whenever both classes had sufficient support; regression random splits were unstratified. The same seed controlled the train-versus-temporary split and the equal calibration-versus-test partition of the temporary set.

For scaffold evaluation, Bemis--Murcko scaffold groups already stored in the processed split files were assigned intact and without target information. Scaffold groups were first shuffled by the prespecified seed, stably ordered by decreasing group size, and allocated to the role with the smallest projected fill ratio relative to its target size. Seed-derived tie ordering was used only when projected ratios were equal. Structural validation required every molecule to receive one role and every scaffold to occur in exactly one role.

Similarity-cluster evaluation used an order-randomized leader algorithm on binary ECFP4 fingerprints. This was not exact Butina clustering. For each seed, molecules were visited in a seeded random order; a molecule joined the most similar existing leader when its Tanimoto similarity was at least 0.60 and otherwise initiated a new cluster. Clusters were then assigned intact to training, calibration, or testing using group size only and the same projected-fill principle. The clustering step was prespecified for endpoints with at most 10,000 molecules; all four core endpoints satisfied this cap. Validation required zero cluster crossover between roles.

\subsection{Frozen predictive models and imbalance regimes}

Four model families were evaluated for each endpoint. Classification used logistic regression, random forest, XGBoost, and a multilayer perceptron on ECFP4. Regression used ridge regression and the corresponding random-forest, XGBoost, and ECFP multilayer-perceptron regressors.

The linear classification pipeline used sparse-compatible standardization and logistic regression with the liblinear solver, 2000 maximum iterations, no class weighting, and the split seed as the random state. Ridge regression used sparse-compatible standardization and $\ell_2$ penalty $\alpha=1.0$. Random forests used 500 trees and the split seed. XGBoost used 300 estimators, maximum depth~4, learning rate~0.05, subsampling~0.90, column subsampling~0.90, histogram tree construction, and the split seed. The multilayer perceptron used sparse-compatible standardization, hidden layers of 256 and 128 units, ReLU activation, $\ell_2$ coefficient $10^{-4}$, a maximum of 300 iterations, early stopping, and the split seed.

Classification was evaluated under two separate regimes. The unweighted regime trained every model on the original training distribution. In the balanced-oversample regime, the minority class was sampled with replacement within the training role until it matched the majority count; the same seeded resampled training matrix was supplied to all four model families. Calibration and test roles were never resampled. Regression required no imbalance regime. Architecture comparisons never mixed weighted and unweighted training as though they were a single model ranking.

\subsection{Predictive performance and probability calibration}

Classification performance was summarized by ROC--AUC, PR--AUC, and balanced accuracy; regression performance used RMSE, MAE, and $R^2$. For binary probability evaluation, predictions were clipped to $[10^{-7},1-10^{-7}]$ before the Brier score and negative log-likelihood were calculated. Calibration was examined in two distinct ten-bin views \citep{guo2017calibration}: positive-probability calibration compared mean $P(Y=1)$ with observed positive frequency, whereas confidence calibration compared maximum class probability with empirical correctness. Expected calibration error was the bin-frequency-weighted absolute gap, and maximum calibration error was the largest occupied-bin gap. These two views were retained separately because calibration of the positive probability is not equivalent to calibration of classification confidence.

\subsection{Marginal and Mondrian conformal prediction}

All headline conformal analyses used $\alpha=0.10$. For binary classification, the least-ambiguous-classification nonconformity score was
\begin{equation}
  s_i = 1-\widehat{p}_{y_i}(x_i),
\end{equation}
where $\widehat{p}_{y_i}(x_i)$ is the probability assigned to the observed class. Marginal split conformal prediction used one calibration distribution for both labels. Its finite-sample order statistic was indexed by
\begin{equation}
  k=\min\!\left\{n_{\mathrm{cal}},\left\lceil(n_{\mathrm{cal}}+1)(1-\alpha)\right\rceil\right\}.
\end{equation}
A candidate label was included when its test nonconformity score did not exceed the corresponding threshold. Mondrian prediction used separate calibration-score distributions and thresholds for class~0 and class~1, providing class-conditional rather than only marginal calibration motivation \citep{cauchois2021knowing,ding2023classconditional}. We reported overall and class-specific coverage, the absolute class-coverage gap, prediction-set size, empty-set rate, singleton rate, and ambiguous two-label rate. These efficiency measures were considered jointly with validity because class-conditional coverage can be repaired by returning weakly informative sets \citep{krstajic2021critical,fisch2022trading}.

For regression, standard split conformal intervals used absolute calibration residuals $|y_i-\widehat{y}_i|$ and the same finite-sample quantile rule. Test intervals were symmetric around the point prediction.

\subsection{Empirical covariate-shift-weighted conformal baseline}

The shift-aware baseline followed the weighted-conformal principle that target-domain coverage can be recovered under covariate shift when the relevant likelihood ratio is known or estimated sufficiently accurately \citep{tibshirani2019covariate}. In each run, calibration ECFP4 vectors were assigned domain label~0 and the unlabeled test ECFP4 vectors domain label~1. Density ratios were estimated with five-fold cross-fitted $\ell_2$ logistic regression using $C=1.0$, liblinear optimization, and no class weighting. Domain-classifier odds were multiplied by $n_{\mathrm{cal}}/n_{\mathrm{test}}$ and clipped to $[0.05,20]$. Weighted conformal quantiles included the test-point mass at infinity. We recorded domain-classifier AUC, effective sample size of the calibration weights, infinite-threshold frequency, empirical coverage, and efficiency.

Because the ratios were estimated from finite molecular samples rather than known analytically, these results are described as empirical shift-aware diagnostics, not as exact guarantees under arbitrary molecular distribution shift. CoDrug provides a recent molecular example of density-aware conformal correction under covariate shift, but uses a different density-estimation strategy \citep{laghuvarapu2023codrug}.

\subsection{Adaptive normalized regression intervals}

The locally adaptive regression baseline divided the calibration role deterministically into equal scale-fitting and conformal-quantile halves using a random generator derived from the split seed. A random-forest scale model with 300 trees, \texttt{min\_samples\_leaf}=20, and \texttt{max\_features}=\texttt{sqrt} was fitted on the first half to predict $\log(1+|y-\widehat{y}|)$. The second half supplied normalized residual scores. The predicted scale was floored at
\begin{equation}
  \max\{0.1\,\mathrm{median}(|r|_{\mathrm{scale}}),\,0.001\},
\end{equation}
and the test half-width was the normalized conformal quantile multiplied by the predicted test scale. Evaluation included coverage, absolute coverage gap, interval width, width coefficient of variation, infinite-interval rate, and Pearson and Spearman association between width and realized absolute error.

\subsection{Applicability-domain diagnostics}

For each test molecule, the primary chemical-domain score was its maximum Tanimoto similarity to the training fingerprints. Mean top-five similarity and unseen-scaffold status were also generated in the pipeline, while the main manuscript emphasizes maximum similarity because it was used consistently in continuous, threshold, and selective analyses. The complementary out-of-domain score was $1-$maximum similarity. We calculated Spearman association between similarity and predictive risk and between similarity and marginal-conformal miscoverage. Out-of-domain AUC measured whether the complementary score discriminated erroneous or miscovered observations.

Prespecified descriptive bins were low domain ($<0.40$), medium domain ($0.40$--$<0.70$), and high domain ($\geq0.70$). Because any fixed structural-similarity threshold is representation and endpoint dependent \citep{kim2025distance}, primary inference used continuous association. Sensitivity analyses swept thresholds from 0.30 to 0.80 and also examined similarity quartiles. Sparse or single-class groups were retained with warnings rather than silently excluded.

\subsection{Selective prediction and retained-class diagnostics}

Classification samples were ranked separately by $1-$maximum class probability and by $1-$maximum Tanimoto similarity; regression used the latter. Retained coverage ranged from 0.10 to 1.00 in increments of 0.05. At every retained fraction, the selected subset was compared with 500 deterministic repeated random-rejection samples of the same size. Classification outputs included ordinary error, balanced error when both classes remained, false-negative and false-positive rates, positive and negative retention, and the shift in retained class balance. Regression outputs included RMSE and MAE.

The area metric was explicitly reported as partial AURC because integration covered retained fractions 0.10--1.00 rather than the full mathematical domain. Selective and random balanced-error curves were integrated only over matched coverage points for which both quantities were finite. This matched-support rule prevents apparent improvements caused solely by dropping different undefined low-coverage points.

\subsection{Aggregation, paired comparisons, and uncertainty}

Metrics were first summarized within a fixed endpoint, split type, seed, model, and imbalance regime. Model-specific means, standard deviations, and two-sided 95\% Student-$t$ intervals were then calculated across confirmatory seeds. Method contrasts were paired within the same endpoint, split, seed, model, and regime. Cross-model headline averages were used only as compact descriptive summaries; model families were not treated as independent inferential replicates. For each paired effect, we additionally counted how many frozen model or model--regime combinations had a model-specific confidence interval entirely above or below zero. No multiplicity-adjusted null-hypothesis testing across model families was used, and null or adverse effects were retained.

\subsection{Stage II: Tox21 cohorts and prospective roles}
\label{sec:methods_tox21}

The design stage used the official Tox21 Challenge cohorts, which comprise nuclear-receptor and stress-response assays measured on environmental chemicals and drugs \citep{huang2016tox21challenge}. The Tox21 10K training collection supplied open labels for model fitting, routing, and internal conformal calibration. A public leaderboard batch of 296 structures supplied architecture-development labels. A separate official file of 647 final structures was the unlabeled target for one independent evaluation. All URLs, archive members, SHA-256 hashes, endpoint order, preprocessing rules, seeds, audit thresholds, promotion gates, and failure actions were recorded in an executable lock before final-label access.

Six primary endpoints---NR-AhR, NR-ER, SR-ARE, SR-ATAD5, SR-MMP, and SR-p53---were selected by the frozen rule of at least 20 positive labels in the public development batch, not by TAME performance. The six remaining official assays were retained as secondary small-sample stress tests and were never used to enlarge the primary denominator. Development used seeds 101--105; independent evaluation used fresh seeds 211--215. Standardized structures overlapping the 10K training cohort were excluded from external domain fitting and endpoint evaluation. An unparseable external structure retained its public identity, received the non-actionable full set $\{0,1\}$, and was excluded with an explicit reason code.

\subsection{Frozen base predictor and ordinary Mondrian set}

For each endpoint and seed, the Tox21 training cohort was partitioned by a stratified, seed-fixed 60/20/20 allocation into fit, router, and internal conformal roles. Four ECFP4 classifiers were fitted on the fit role: logistic regression, random forest, extra trees, and Bernoulli naive Bayes. The tree ensembles used 120 estimators and square-root feature sampling; the two linear-probabilistic models and all remaining hyperparameters were fixed in the executable lock. An $\ell_2$ logistic router was trained on the four component logits in the router role.

For routed probability $p(x)$, the candidate-label nonconformity scores were $s_0(x)=p(x)$ and $s_1(x)=1-p(x)$. Ordinary class-conditional Mondrian thresholds at $\alpha_0=\alpha_1=0.10$ formed the protected baseline $S_M(x)$. Within class $y$, the finite-sample rank was
\begin{equation}
 k_y=\left\lceil(n_y+1)(1-\alpha_y)\right\rceil ,
\end{equation}
with an unavailable rank mapped to positive infinity. This base layer was fixed before either external cohort was evaluated.

\subsection{Label-free transport views and audit certificate}
\label{sec:methods_transport}

TAME estimated source-to-target density ratios from two complementary views that used external structures and model outputs but no external labels. The physicochemical view contained molecular weight, logP, topological polar surface area, hydrogen-bond donor and acceptor counts, rotatable-bond count, ring count, fraction sp$^3$, and heavy-atom count. The score view contained the four component logits and the routed logit. A 2048-bit ECFP domain classifier was retained as a mandatory diagnostic comparator but was not eligible to enter the envelope after failing its development certificate.

Each view used five-fold cross-fitted $\ell_2$ logistic domain classification for the domain-AUC audit and the same regularized estimator refitted on all unlabeled source and target covariates for the working density ratio. Posterior odds were prior-corrected by $n_{\mathrm{source}}/n_{\mathrm{target}}$ and clipped to $[0.05,20]$. The weighted source scores and query mass at positive infinity then defined candidate-specific weighted Mondrian thresholds \citep{tibshirani2019covariate,laghuvarapu2023codrug}. Unlike KMM-CP, which estimates covariate-shift weights through constrained kernel mean matching and adds a support-overlap selection step \citep{laghuvarapu2026kmmcp}, TAME retained the frozen logistic ratio estimator and used the complete certificate below as an activation gate. KMM-CP was not part of the locked comparator suite; adding a new final-label comparison after unsealing would violate the evidence boundary, so a head-to-head test requires a new prospectively locked study.

A view was active only if all six prespecified checks passed: (i) at least 100 source and 100 target structures; (ii) at least 20 internal conformal examples in each true class; (iii) total and per-class effective sample size at least 25\% of the corresponding raw count; (iv) no more than 30\% of weights at either clipping boundary; (v) cross-fitted domain AUC no greater than 0.99; and (vi) the weighted standardized mean gap no more than 1.05 times the unweighted gap. These diagnostics are empirical transport certificates. They do not prove that estimated ratios are correct or restore exact coverage under arbitrary molecular distribution shift.

\subsection{TAME envelope, fallback, and structural invariants}
\label{sec:methods_tame}

Public development showed a mean coverage deficit for label~0 greater than two percentage points and no mean deficit for label~1. The protected label was therefore frozen as label~0 before the final EPA evaluation. This label choice used public development outcomes; only the application of the two transport views to the final target was label free.

Let $\bar S_M(x)=S_M(x)$ when the baseline is nonempty and $\bar S_M(x)=\{0,1\}$ otherwise. Let $S_{\mathrm{phys}}(x)$ and $S_{\mathrm{score}}(x)$ be the two weighted Mondrian sets. If both views pass their complete audits,
\begin{equation}
 S_{\mathrm{TAME}}(x)
 =\bar S_M(x)\cup
 \left\{0:\,0\in S_{\mathrm{phys}}(x)\cap S_{\mathrm{score}}(x)\right\}.
 \label{eq:tame}
\end{equation}
If either approved view fails, every transport augmentation is disabled and $S_{\mathrm{TAME}}(x)=\bar S_M(x)$. Thus a failed audit is observable and produces an exact ordinary-baseline fallback rather than a silently relaxed certificate.

Equation~\ref{eq:tame} yields three target-label-independent invariants: the TAME set contains the ordinary Mondrian set, TAME emits no empty set, and every TAME singleton was already a baseline singleton. Consequently, empirical coverage for either class cannot decrease on the same rows, and wrong-singleton exposure cannot increase. TAME can convert a baseline singleton into an explicit two-label defer, but cannot create a new confident singleton. Correct-singleton yield can therefore decrease but cannot increase; this motivated a prespecified efficiency guardrail rather than an efficiency-superiority claim.

\subsection{Development gate and prediction-to-label firewall}
\label{sec:methods_firewall}

The public development denominator was fixed at six primary endpoints by five seeds (30 endpoint--seed cells). Promotion required all cells to complete, zero baseline-inclusion violations, zero TAME empty sets, zero newly created singletons, nonnegative mean coverage change for both classes, nonpositive wrong-singleton-exposure change for both classes, mean MacroCSY change of at least $-0.05$, and both approved views active in at least 80\% of cells. Failure of any check would stop execution before final-label download or parsing, and the failed gate would remain recorded.

After the gate passed, fresh-seed models generated every comparator and TAME prediction from the 647 final structures. The sealed prediction file, final transport-audit table, and executable lock were hashed into a promotion record with \texttt{final\_labels\_opened=false}. Only then could the locked final-label bytes be acquired, their hash verified, and labels identity-joined for a single evaluation. The parser had no callable path around the promotion record, and contract tests exercised missing or invalid records. No endpoint, seed, transport view, threshold, protected label, or margin could be changed after final-label inspection.

\subsection{Primary estimand, efficiency metric, and inference}
\label{sec:methods_inference}

For class $y$, coverage was the fraction of true-$y$ compounds whose set contained $y$. The primary endpoint-level quantity was the smaller of the two class coverages. Correct-singleton yield for class $y$ was the fraction of true-$y$ compounds receiving the singleton $\{y\}$; MacroCSY was the unweighted mean of the two class-specific yields. Wrong-singleton exposure was analogously the fraction receiving the singleton containing the incorrect label. Ambiguous and empty rates completed the descriptive set metrics.

The primary estimand was the endpoint--seed-equal mean change in minimum-class coverage for TAME versus the ordinary routed Mondrian baseline across the six primary endpoints and five independent seeds. Its 95\% interval used a frozen hierarchical bootstrap: endpoints were resampled first; within every drawn endpoint, labeled final compounds were then resampled separately within true class, preserving class counts; all five seeds were retained. The procedure used 2000 draws and seed~44021. Secondary estimands included MacroCSY change, class-specific coverage, wrong-singleton exposure, ambiguity, and empty-set rate. Interpretation was hierarchical: a positive primary point estimate had to be accompanied by a MacroCSY point estimate no worse than the frozen $-5$-point guardrail. On the class-balanced MacroCSY scale, this threshold permits at most five fewer correct singleton outputs per 100 cases. It was fixed after public development as an architecture-selection tolerance and before final-label access; it was not calibrated to clinical or regulatory utility. The executable lock names the field a non-inferiority margin, but no confidence-interval-based non-inferiority test for MacroCSY was prespecified. We therefore treat it as a point-estimate efficiency guardrail and make no formal non-inferiority claim. No multiplicity-adjusted secondary superiority claim was authorized automatically.

A post-evaluation implementation audit found that the first bootstrap report traversed unique resampled endpoints through an unordered Python set. This affected only how a fixed random stream was assigned within the percentile calculation. A hash-verified inference-only repair preserved first-occurrence endpoint order and recomputed the interval without refitting a model, regenerating a prediction, changing a label, or altering either point estimate. The repaired deterministic interval is the publication-facing value reported here; the original sealed report remains versioned in the audit chain.
